% ============================================================
%  chapters/ch24-meaning-as-geometry.tex
% ============================================================

\chapter{Meaning as Geometry}

\chapterepigraph{%
  Meaning is not what a symbol points to.\\
  Meaning is the shape of the region\\
  from which the symbol cannot escape.
}{Flyxion, \textit{Semantic Infrastructure}}

\sidenote{Connects to\\App.\,L §L.23\\and the RSVP\\semantic field\\equations}

The previous two chapters developed the agent's internal world as a
latent simulation (Chapter~22) and the temporal structure of that
simulation as a memory chain (Chapter~23). This chapter asks: what is
\emph{meaning} in this framework?

The answer is geometric. Meaning is not a correspondence relation
between symbols and objects (the referential theory), nor is it a
use-pattern (the Wittgensteinian theory), nor is it a mental representation
(the cognitive theory). It is a geometric structure in the accessibility
landscape: the shape of the constraint basin that a semantic unit inhabits.

\section{Semantic Fields}

In the RSVP framework, meaning lives in a semantic manifold $M_s$ —
the space of possible semantic states — equipped with an accessibility
entropy field $S_s : M_s \to \mathbb{R}$ (Appendix~L).

\begin{definition}{Semantic Field}{semantic-field}
  The \emph{semantic field} of a symbol $\sigma$ is the region of the
  semantic manifold $M_s$ consistent with $\sigma$:
  \[
    \mathcal{F}(\sigma) = \{m \in M_s : \Phi_s(\sigma, m) \geq \tau\}
  \]
  where $\Phi_s : \Sigma \times M_s \to [0,1]$ is the compatibility
  function (how well semantic state $m$ is expressed by symbol $\sigma$)
  and $\tau$ is a compatibility threshold.
\end{definition}

The semantic field is not a point but a region. The word ``tree'' has
a semantic field that includes all semantic states involving tree-like
things: oaks, concepts of trees, tree-shaped graphs, decision trees,
family trees. The boundary of this field — where $\Phi_s(\sigma, m) = \tau$
— is the edge of the symbol's meaning.

\section{Qualia as Invariants}

The hardest problem of consciousness is the problem of \emph{qualia}:
why does red look like red? Why does pain feel like pain? Why is there
``something it is like'' to experience these things?

The geometric framework offers a partial answer. Qualia are not
properties of the environment or of symbols — they are \emph{invariants
of the agent's projection structure}.

\begin{definition}{Quale as Projection Invariant}{quale-invariant}
  A \emph{quale} associated with stimulus class $S$ is a property of
  the agent's projection $\Pi_A$ that is invariant across all stimuli
  in $S$: it is preserved by $\Pi_A$ for all $e \in S$ and varies
  across different stimulus classes.
  Formally, it is an element of the fiber $\ker(\Pi_A|_S)^\perp$ —
  the component of the projection that distinguishes $S$ from other
  classes.
\end{definition}

Red looks like red because the quale of redness is an invariant of the
human visual projection $\Pi_{\text{human}}$ for wavelengths near 700nm.
It is not a property of the wavelength itself — other projections (other
species' visual systems) have different invariants for the same wavelength.
The quale is the agent-specific structure that the projection extracts.

\begin{remark}
  This does not solve the hard problem — it relocates it. We still do
  not know why projection invariants have subjective character at all.
  But it does explain the \emph{structure} of qualia: why they are
  stable (invariance under stimulus variation), agent-specific (depend
  on $\Pi_A$), and incommunicable (fibers of projections are invisible
  to other agents with different projections).
\end{remark}

\section{Constraint Basins}

The most powerful geometric concept for meaning is the \emph{constraint
basin}: the region of the accessibility landscape that a semantic state
is drawn toward by the accessibility gradient flow.

\begin{definition}{Constraint Basin}{constraint-basin}
  The \emph{constraint basin} of a semantic attractor $m^* \in M_s$
  is the set of initial states from which the accessibility gradient
  flow converges to $m^*$:
  \[
    \mathcal{B}(m^*) = \{m \in M_s : \lim_{t\to\infty}
    \phi_t(m) = m^*\}
  \]
  where $\phi_t$ is the flow of the accessibility gradient vector field
  $\nabla S_s$.
\end{definition}

Words and concepts are attractors in the semantic manifold; their meanings
are their constraint basins. The basin of the word ``tree'' is the set of
all semantic states that, under the pressure of constraint satisfaction,
converge to the tree-concept attractor. Words with overlapping basins
are synonyms; words with nearby but non-overlapping basins are near-synonyms.

\begin{figure}[h]
  \centering
  \begin{tikzpicture}[scale=1.2]
    % semantic-basins.tex — constraint basins in semantic manifold
\tikzset{
  attractor/.style={circle, fill=RSVPdeep, minimum size=7pt, inner sep=0},
  basin/.style={fill opacity=0.2, draw opacity=0.8, thick},
}
  % Basin 1: "tree" concept (left)
  \fill[RSVPmid, basin] (-2.2, 0) ellipse (1.3cm and 0.85cm);
  \draw[RSVPmid, thick] (-2.2, 0) ellipse (1.3cm and 0.85cm);
  \node[attractor] at (-2.2, 0) {};
  \node[RSVPdeep, font=\small\itshape, below=6pt] at (-2.2, 0) {tree};

  % Basin 2: "forest" concept (center-left, overlapping slightly)
  \fill[RSVPaccent, basin] (-0.5, 0.3) ellipse (1.1cm and 0.7cm);
  \draw[RSVPaccent, thick] (-0.5, 0.3) ellipse (1.1cm and 0.7cm);
  \node[attractor, fill=RSVPaccent] at (-0.5, 0.3) {};
  \node[RSVPaccent!80, font=\small\itshape, below=6pt] at (-0.5, 0.3) {forest};

  % Basin 3: "network" concept (right)
  \fill[RSVPpurple, basin] (1.8, -0.1) ellipse (1.2cm and 0.75cm);
  \draw[RSVPpurple, thick] (1.8, -0.1) ellipse (1.2cm and 0.75cm);
  \node[attractor, fill=RSVPpurple] at (1.8, -0.1) {};
  \node[RSVPpurple, font=\small\itshape, below=6pt] at (1.8, -0.1) {network};

  % Gradient flow arrows (inward toward attractors)
  % toward tree
  \draw[-{Stealth[length=3.5pt]}, RSVPmid!70, thin] (-3.2, 0.5) -- (-2.65, 0.2);
  \draw[-{Stealth[length=3.5pt]}, RSVPmid!70, thin] (-3.0, -0.6) -- (-2.65, -0.2);
  \draw[-{Stealth[length=3.5pt]}, RSVPmid!70, thin] (-1.5, 0.7) -- (-2.0, 0.3);
  % toward forest
  \draw[-{Stealth[length=3.5pt]}, RSVPaccent!70, thin] (-0.5, 1.5) -- (-0.5, 0.7);
  \draw[-{Stealth[length=3.5pt]}, RSVPaccent!70, thin] (0.5, 1.0) -- (0.1, 0.6);
  % toward network
  \draw[-{Stealth[length=3.5pt]}, RSVPpurple!70, thin] (2.8, 0.5) -- (2.25, 0.15);
  \draw[-{Stealth[length=3.5pt]}, RSVPpurple!70, thin] (2.6, -0.8) -- (2.2, -0.35);

  % Basin boundary label
  \node[RSVPdeep!60, font=\footnotesize\itshape] at (0.6, -1.3)
    {constraint basin boundaries};
  \node[RSVPmid!80, font=\footnotesize\itshape] at (-3.0, 1.2)
    {$\nabla S_s$};

  \end{tikzpicture}
  \caption{Constraint basins in the semantic manifold. Each basin
    (shaded region) surrounds an attractor (filled circle) corresponding
    to a concept. The accessibility gradient $\nabla S_s$ drives semantic
    states toward their nearest attractor. Basin boundaries (thick curves)
    are the separatrices of the gradient flow — the edges of meaning.}
  \label{fig:semantic-basins}
\end{figure}

\section{Coherence Structures}

\begin{tcolorbox}[enhanced, breakable,
  colback=RSVPpurple!8, colframe=RSVPpurple, arc=4pt, boxrule=0.6pt,
  title={\bfseries\sffamily Technical Note: Whitney Stratification and Semantic Phase Transitions},
  coltitle=white, attach boxed title to top left={yshift=-2mm,xshift=6mm},
  boxed title style={colback=RSVPpurple,sharp corners},
  top=4mm,left=4mm,right=4mm,bottom=4mm]

The semantic manifold $M_s$ is not a smooth uniform space — it is
\emph{stratified}: divided into zones (strata) where qualitatively
different cognitive rules apply. Mathematical reasoning, spatial reasoning,
emotional reasoning, and social reasoning occupy different strata.

\textbf{Whitney's condition~B} governs crossing between strata.
When a cognitive trajectory approaches the boundary between two strata,
it must approach \emph{tangentially} — gliding along the edge, allowing
the semantic structures of one regime to continuously map onto the other.
A tangential crossing preserves coherence.

A \emph{transversal} crossing — approaching the boundary head-on — produces
a \emph{non-admissible phase transition}: catastrophic loss of semantic
structure. In human terms: suddenly losing your train of thought when
switching between two very different modes of reasoning. In AI terms:
an incoherent output produced when the system jumps between logical strata
without preserving the underlying geometry.

The HYDRA TARTAN tiling architecture addresses this by decomposing the
manifold into locally uniform tiles with \emph{annotated noise} at each
boundary — an explicit measurement of residual variability that allows the
system to recognize when it is near a stratum boundary and slow down
accordingly, approaching tangentially rather than transversally.

Practical consequence: the human experience of ``shifting gears'' between
different types of thinking — from formal calculation to creative
imagination — has a precise mathematical form. The discomfort or confusion
of the gear-shift is the cognitive cost of approaching a stratum boundary
non-tangentially. Good thinking involves learning the topography of one's
own semantic manifold well enough to always cross boundaries tangentially.
\end{tcolorbox}


Meaning is not just the identity of individual concepts but the
\emph{coherence structure} that holds multiple concepts together.

\begin{definition}{Semantic Coherence Field}{sem-coherence}
  The \emph{semantic coherence field} $C_s : M_s^n \to [0,1]$ measures
  how well $n$ semantic states mutually support each other — how much
  each state's activation reduces the constraint violation of the others.
  High coherence corresponds to a cluster of mutually reinforcing concepts;
  low coherence corresponds to a collection of unrelated or contradictory ones.
\end{definition}

A text has high coherence when its constituent concepts occupy a
high-coherence region of $M_s^n$. Understanding a text is navigating
to a region of high coherence. Confusion is occupying a low-coherence
region and not finding a path to a higher one.

This explains why metaphor works. A good metaphor maps two domains
— whose concepts usually occupy distant regions of $M_s$ — onto a
high-coherence configuration that illuminates both. ``An argument is
a building'' is coherent because the structure-stability-foundation
cluster of the building domain maps onto a high-coherence configuration
with the claim-support-premises cluster of the argument domain.

\section{Zero Ontology and Semantic Geometry}

The zero ontology of Pearce (discussed in the preface to Part~VI)
proposes that the universe has zero net information content. The
semantic geometry framework offers an interesting perspective on this.

If the universe is an accessibility landscape and meaning is the
geometry of constraint basins, then the zero-information universe
is one in which all constraint basins have zero depth — no concept
has more gravitational pull than any other. This is not a universe
of rich meaning but a universe of undifferentiated accessibility:
maximum $S_s$ everywhere, minimum $\nabla S_s$.

The formation of meaning — the deepening of constraint basins — is
therefore equivalent to the decrease of global accessibility entropy:
meaning is purchased at the cost of future possibility. A highly
meaningful conceptual structure is a highly constrained one.

\begin{namedthm}{The Meaning–Accessibility Tradeoff}
  In any accessibility landscape, the emergence of stable semantic
  attractors (concepts, meanings, values) corresponds to a decrease
  in global accessibility entropy. There is a fundamental tradeoff:
  \[
    \text{Meaning depth} \;\propto\; -\Delta S_{\text{global}}.
  \]
  A universe with zero semantic content has maximum accessibility;
  a universe with rich semantic structure has consumed much of its
  accessibility in forming constraint basins.
\end{namedthm}

\begin{exercise}
  Show that two words $\sigma_1, \sigma_2$ are synonyms in the
  semantic field framework iff $\mathcal{F}(\sigma_1) \cap
  \mathcal{F}(\sigma_2) \neq \emptyset$ and neither is a proper
  subset of the other. What does a word being a proper subset of
  another correspond to semantically?
\end{exercise}

\begin{exercise}
  The constraint basin $\mathcal{B}(m^*)$ depends on the accessibility
  gradient flow, which depends on the accessibility field $S_s$.
  Different agents with different constraint structures will have
  different flows and therefore different basins for the same attractor.
  Formalize the claim that ``the same word means different things
  to different people'' in these terms.
\end{exercise}

\begin{exercise}[title={Connecting frameworks}]
  Pearce argues that ``nothing is difficult to define'' because any
  specification of nothing already contains structure. Formalize this
  in the constraint-basin framework: what is the constraint basin of
  the concept ``nothing''? Is it empty, unbounded, or something else?
  What does its shape tell us about why nothingness is philosophically
  problematic?
\end{exercise}
